\documentclass[10pt]{article}
\usepackage{precalculus-article}
\usepackage{precalculus-boxes}
\usepackage{graphicx}
\graphicspath{{images/}}
\newcommand{\TallMath}[1]{$\displaystyle #1\rule[-1.4em]{0pt}{3.2em}$}
\newcommand{\CourseName}{AP Precalculus}
\newcommand{\SchoolYear}{2026--2027}
\newcommand{\MeetingLength}{55 minutes}
\newcommand{\UnitNumberName}{Unit 4: Functions Involving Parameters, Vectors, and Matrices \quad}
\newcommand{\LessonNumberName}{Lesson 4.5: Implicitly Defined Functions}

\begin{document}
\begin{center}
    {\Large\bfseries \CourseName: \SchoolYear \\
    {\small\bfseries \UnitNumberName \LessonNumberName}}
\end{center}

\begin{tcolorbox}[colback=sky, colframe=navy, boxrule=0.9pt, arc=2mm,
    left=2mm, right=2mm, top=1.5mm, bottom=1.5mm]
    \textbf{Primary Objective:} Students will graph equations in two variables by finding
    solutions and identifying one or more implicit functions, determine how the variables
    co-vary by examining the sign of their ratio of change, and locate where the rate of
    change is zero or undefined and interpret this geometrically. (Big Idea: Functions)
\end{tcolorbox}

\renewcommand{\arraystretch}{1.6}
\begin{skillbox}[Priority Ideas \& Skills]{goldbox}
\begin{tabularx}{\linewidth}{@{}X X@{}}
\begin{minipage}[t]{0.48\linewidth}\vspace{0pt}
    \textbf{AP Skill 2.B} --- Construct equivalent representations\\[4pt]
    \textbf{AP Skill 3.A} --- Describe characteristics of a function\\[6pt]
    \textbf{Learning Objectives:}
    \begin{itemize}[leftmargin=1.2em, itemsep=2pt]
        \item \textbf{LO 4.5.A:} Construct a graph of an equation involving two variables.
        \item \textbf{LO 4.5.B:} Determine how the two quantities related in an implicitly
              defined function vary together.
    \end{itemize}
\end{minipage}
&
\begin{minipage}[t]{0.48\linewidth}\vspace{0pt}
    \textbf{Key Understandings (EKs):}
    \begin{itemize}[leftmargin=1.2em, itemsep=2pt]
        \item An equation in two variables can implicitly describe one or more functions.
              (EK 4.5.A.1)
        \item A graph can be built by substituting values and solving for the other
              variable. (EK 4.5.A.2)
        \item Solving for one variable yields an explicit function whose graph is part or
              all of the equation's graph. (EK 4.5.A.3)
        \item If $\Delta y/\Delta x$ is positive, the variables change in the same
              direction; if negative, they change in opposite directions. (EK 4.5.B.1)
        \item $\Delta y/\Delta x = 0$ signals a horizontal interval; undefined signals a
              vertical interval. (EK 4.5.B.2)
    \end{itemize}
\end{minipage}
\end{tabularx}
\end{skillbox}

\begin{skillbox}[{Vocabulary, Concepts, \& Theorems}]{greenbox}
\renewcommand{\arraystretch}{1.5}
\begin{tabularx}{\linewidth}{@{} l X @{}}
\textbf{Implicit equation} & An equation in two variables where neither variable is isolated;
  it may describe one or more functions. \\
\textbf{Explicitly defined function} & A function of the form $y = f(x)$, where $y$ is
  solved for in terms of $x$. \\
\textbf{Co-variation} & The way two quantities change together; characterized by the sign
  and magnitude of $\Delta y/\Delta x$. \\
\textbf{Rate of change} & $\Delta y/\Delta x$: the ratio of the change in $y$ to the change
  in $x$ between two nearby points. Positive $\Rightarrow$ same direction;
  negative $\Rightarrow$ opposite directions. \\
\end{tabularx}
\end{skillbox}

\begin{skillbox}[Activate Prior Knowledge \& Spiral Review]{sky}
    \begin{tabularx}{\linewidth}{@{}X c@{}}
        \begin{minipage}[t]{0.55\linewidth}\vspace{0pt}
            \textbf{Students should already be able to:}
            \begin{itemize}[leftmargin=1.2em, itemsep=2pt]
                \item Solve a two-variable equation for one variable algebraically.
                \item Evaluate $y = f(x)$ for specific $x$ values and plot $(x, y)$ pairs.
                \item Compute $\Delta y/\Delta x$ from a table of values.
                \item Use the square-root property ($y^2 = k \Rightarrow y = \pm\sqrt{k}$).
            \end{itemize}
            \vspace{4pt}\textit{Warm-Up covers:} Solving $x^2+y^2=100$ for $y$;
            reading a ratio from a table; finding $y$-values for $y^2 = x+4$.
        \end{minipage}
        &
        \begin{minipage}[t]{0.38\linewidth}\vspace{0pt}\centering
            % Authored warm-up; thumbnail not embedded (compiles to target/).
        \end{minipage}
    \end{tabularx}
\end{skillbox}

\begin{skillbox}[Hook]{sky}
    \textit{``The equation $x^2 + y^2 = 25$ is not a function --- but it describes a circle.
    Can we still talk about how $x$ and $y$ change together? Can we split it into functions?''}\\[4pt]
    Ask students: Is $x^2 + y^2 = 25$ a function? Why or why not? (No --- fails the vertical-line
    test.) What if we only look at the top half? Guide them to see that each semicircle \emph{is}
    a function.
\end{skillbox}

\begin{skillbox}[Lesson --- Day 1: Graphing Implicit Equations]{sky}
\begin{multicols}{2}
\textbf{Step 1 --- What is an implicit equation?}\\
An equation like $x^2 + y^2 = 25$ involves two variables but neither is isolated.
We call it \textbf{implicit} because functions are hidden inside it.

\textbf{Pose:} ``How could we graph this?'' \\
\textbf{Guide:} Substitute specific $x$ values and solve for $y$.

\textbf{Example:} $x = 3$ in $x^2+y^2=25$:
\[ 9 + y^2 = 25 \implies y^2 = 16 \implies y = \pm 4 \]
Two points, $(3, 4)$ and $(3, -4)$, for one $x$-value.

\textbf{Step 2 --- Identifying the explicit pieces.}\\
Solve $x^2 + y^2 = 25$ for $y$:
\[ y = \pm\sqrt{25-x^2} \]
Two explicit functions on $[-5, 5]$:
\begin{itemize}[leftmargin=1em, itemsep=1pt]
    \item $f_1(x) = \sqrt{25-x^2}$ \;(upper semicircle)
    \item $f_2(x) = -\sqrt{25-x^2}$ \;(lower semicircle)
\end{itemize}

\columnbreak

\textbf{Step 3 --- Table for $x^2+y^2=25$.}\\[2pt]
\renewcommand{\arraystretch}{1.3}
\begin{tabular}{c|c|c}
$x$ & $y$ (upper) & $y$ (lower) \\\hline
$-5$ & $0$  & $0$  \\
$-3$ & $4$  & $-4$ \\
$0$  & $5$  & $-5$ \\
$3$  & $4$  & $-4$ \\
$5$  & $0$  & $0$  \\
\end{tabular}\\[6pt]
Have students verify each row by substitution.

\textbf{Step 4 --- Second example: $xy = 4$.}\\
Solve: $y = 4/x$.  One explicit function (two branches, no repeated $y$).\\[4pt]
\begin{tabular}{c|c}
$x$ & $y$ \\\hline
$1$ & $4$  \\ $2$ & $2$ \\ $4$ & $1$ \\
$-1$ & $-4$ \\ $-2$ & $-2$ \\
\end{tabular}\\[4pt]
\textbf{Pose:} ``As $x$ increases from 1 to 4, what happens to $y$?''
(Decreases from 4 to 1.)
\end{multicols}
\end{skillbox}

\begin{skillbox}[Lesson (cont.) --- Day 2: Co-variation]{sky}
\begin{multicols}{2}
\textbf{Step 5 --- Co-variation sign.}\\
Pick two nearby points on the upper semicircle of $x^2+y^2=25$:
$(3,\,4)$ and $(4,\,3)$.
\[ \frac{\Delta y}{\Delta x} = \frac{3-4}{4-3} = -1 \]
Negative $\Rightarrow$ as $x$ increases, $y$ decreases. \textbf{Opposite directions.}

\textbf{Pose:} ``What would a \emph{positive} ratio mean?''\\
(Both variables increasing, or both decreasing together.)

\textbf{Step 6 --- Where is the ratio zero or undefined?}\\
\begin{itemize}[leftmargin=1em, itemsep=2pt]
    \item Near $(0,5)$: $\Delta y/\Delta x \to 0$.
          $y$ barely changes --- \textbf{horizontal} behavior at the top.
    \item Near $(5,0)$: $\Delta y/\Delta x$ becomes very large (undefined limit)
          --- \textbf{vertical} behavior at the right end.
\end{itemize}

\columnbreak

\textbf{Summary:}\\[2pt]
\renewcommand{\arraystretch}{1.4}
\begin{tabularx}{\columnwidth}{@{}X l@{}}
\textbf{Ratio} & \textbf{Interpretation} \\\hline
Positive & Same direction (both $\uparrow$ or both $\downarrow$) \\
Negative & Opposite directions \\
Zero & $y$ momentarily constant (horizontal) \\
Undefined & $x$ momentarily constant (vertical) \\
\end{tabularx}

\vspace{6pt}
\textbf{$xy = 4$ example:}\\
Points $(2,2)$ and $(3,\,4/3)$:
\[ \frac{\Delta y}{\Delta x} = \frac{4/3 - 2}{3-2} = -\tfrac{2}{3} \]
Negative: as $x$ increases, $y$ decreases. Always true when $xy = k > 0$.
\end{multicols}
\end{skillbox}

\begin{skillbox}[Active Monitoring]{redbox}
\textbf{Circulate and check for:}
\begin{itemize}[leftmargin=1.5em, itemsep=2pt]
    \item Students who find only one $y$-value when two exist (e.g., forgetting $y = -4$ when
          $y^2 = 16$). Prompt: ``$y^2 = 16$ --- what are ALL values of $y$?''
    \item Students who compute $\Delta x/\Delta y$ instead of $\Delta y/\Delta x$.
    \item Students who confuse the sign of the ratio with the sign of the $y$-values themselves.
    \item Students who cannot identify which piece of the circle is a function; prompt with the
          vertical-line test.
\end{itemize}
\textbf{Cold-call prompts:}
\begin{itemize}[leftmargin=1.5em, itemsep=2pt]
    \item ``If $\Delta y/\Delta x$ is negative, what does that tell us about co-variation?''
    \item ``How many explicit functions does $x^2 + y^2 = 25$ define? Name them.''
    \item ``At which points on the circle is the rate of change equal to zero?''
\end{itemize}
\end{skillbox}

\begin{skillbox}[Group Work \& Differentiation]{redbox}
\begin{multicols}{3}
\textbf{Tier R --- Remediate}\\
For $x^2 + y^2 = 16$, find four solutions by substituting $x = -4,\; 0,\; 4$ and one other
value. Determine whether $y$ is increasing or decreasing as $x$ increases from $-4$ to $0$
on the upper half.

\columnbreak
\textbf{Tier A --- Approaching Proficiency}\\
For $x^2 + 4y^2 = 16$, find solutions for $x = -4,\,-2,\,0,\,2,\,4$, identify the two
explicit functions, and determine where $\Delta y/\Delta x$ changes sign.

\columnbreak
\textbf{Tier E --- Extension}\\
For $x^2 - y^2 = 9$, identify two explicit functions and state their domains. Describe
co-variation in each quadrant. What happens at $x = 3$ and $x = -3$?
\end{multicols}
\end{skillbox}

\begin{skillbox}[Individual Work \& Assessment]{redbox}
\textbf{Exit Ticket (2 items):}
\begin{enumerate}[leftmargin=1.5em, itemsep=4pt]
    \item The equation $4x^2 + 9y^2 = 36$ (an ellipse). For $x = 3$, find the $y$-values.
          For $x = 0$, find the $y$-values.
    \item Using $(3,0)$ and $(0,2)$ from $4x^2 + 9y^2 = 36$, determine whether co-variation
          is positive or negative as we move from $(3,0)$ toward $(0,2)$ along the upper half
          of the ellipse.
\end{enumerate}
\textbf{Assessment note:} Sort exit tickets by whether students (a) found both $y$-values
and (b) correctly identified the sign of co-variation. Use results to form Tier R/A/E groups
for Day 2 re-teach if needed.
\end{skillbox}

\begin{skillbox}[Reinforcement \& Extension]{goldbox}
\textbf{Homework:} 5--6 problems including solving for $y$ from an implicit equation, finding
ordered-pair solutions, identifying explicit function pieces, determining co-variation sign,
a contextual budget-constraint problem, and one AP-style MC item.\\[4pt]
\textbf{Extension:} Explore $x^2 - y^2 = 9$ (a hyperbola). Identify its two branches,
describe co-variation in each branch, and find where the rate of change is zero or
undefined.\\[4pt]
\textbf{Preview of Lesson 4.6:} We will extend implicit equations to parametric
representations, where both $x$ and $y$ are expressed as functions of a third variable $t$.
\end{skillbox}

\end{document}
